\clearpage\pagestyle{otMain}\fontsize{11}{15}\selectfont\corrigesChapitre{15}{Transformations et homothéties}
\begin{blocLivre}{cor-15-1}\begin{corrige}{1}{Une symétrie axiale}\demarche{Reporter perpendiculairement de l’autre côté de l’axe.}
On garde l’ordonnée et on change le signe de l’abscisse : $\res{A'(-2;1),B'(-4;1),C'(-3;3)}$. L’axe est la médiatrice des trois segments point-image.
\label{fig-cor-15-1}\figureGeo[corrige]{transfo-01}
\end{corrige}\end{blocLivre}
\begin{blocLivre}{cor-15-2}\begin{corrige}{2}{Une symétrie centrale}\demarche{Placer le centre au milieu du segment.}
O est le milieu de [AA’], donc $\res{A'(-3;-2)}$. Les deux coordonnées changent de signe.
\label{fig-cor-15-2}\figureGeo[corrige]{transfo-02}
\end{corrige}\end{blocLivre}
\begin{blocLivre}{cor-15-3}\begin{corrige}{3}{Une translation}\demarche{Répéter le même déplacement pour tous les sommets.}
Chaque point se déplace de trois unités vers la droite et d’une vers le haut. $\res{A'(4;3),B'(6;3),C'(5;5)}$. Les côtés gardent longueurs et directions.
\label{fig-cor-15-3}\figureGeo[corrige]{transfo-03}
\end{corrige}\end{blocLivre}
\begin{blocLivre}{cor-15-4}\begin{corrige}{4}{Un quart de tour}\demarche{Respecter le sens et conserver la distance au centre.}
Le quart de tour donne $\res{A'(-1;3)}$. On vérifie $OA^2=10=OA'^2$ et l’angle de 90 degrés dans le sens demandé.
\label{fig-cor-15-4}\figureGeo[corrige]{transfo-04}
\end{corrige}\end{blocLivre}
\begin{blocLivre}{cor-15-5}\begin{corrige}{5}{Longueurs conservées}\demarche{Utiliser l’invariance plutôt que mesurer le dessin.}
Une rotation conserve chaque longueur. Le périmètre reste $4+6+7=\res{17\text{ cm}}$, indépendamment de l’angle de rotation.
\label{fig-cor-15-5}\figureGeo[corrige]{transfo-05}
\end{corrige}\end{blocLivre}
\begin{blocLivre}{cor-15-6}\begin{corrige}{6}{Rapport positif}\demarche{Doubler les distances au centre sans changer de côté.}
Chaque coordonnée est doublée : $\res{A'(4;2),B'(6;0),C'(2;6)}$. Les images restent sur les mêmes demi-droites issues de O.
\label{fig-cor-15-6}\figureGeo[corrige]{transfo-06}
\end{corrige}\end{blocLivre}
\begin{blocLivre}{cor-15-7}\begin{corrige}{7}{Rapport négatif}\demarche{Le signe règle le côté ; la valeur absolue règle la longueur.}
On multiplie les coordonnées par $-1{,}5$ : $\res{A'(-3;-1{,}5),B'(-4{,}5;0),C'(-1{,}5;-4{,}5)}$. Les images sont de l’autre côté du centre, à 1,5 fois la distance.
\label{fig-cor-15-7}\figureGeo[corrige]{transfo-07}
\end{corrige}\end{blocLivre}
\begin{blocLivre}{cor-15-8}\begin{corrige}{8}{Retrouver le rapport}\demarche{Distinguer rapport signé et longueur positive.}
La valeur absolue est $10/4=2{,}5$. Les points sont de côtés opposés : $\res{k=-2{,}5}$. Une distance OA’ n’est pourtant jamais négative.
\label{fig-cor-15-8}\figureGeo[corrige]{transfo-08}
\end{corrige}\end{blocLivre}
\begin{blocLivre}{cor-15-9}\begin{corrige}{9}{Périmètre et aire}\demarche{Appliquer le facteur adapté à la dimension.}
Les longueurs sont multipliées par 3 : périmètre $\res{54\text{ cm}}$. L’aire est multipliée par $(-3)^2=9$ : $\res{180\text{ cm}^2}$. Aucun résultat de mesure n’est négatif.
\label{fig-cor-15-9}\figureGeo[corrige]{transfo-09}
\end{corrige}\end{blocLivre}
\begin{blocLivre}{cor-15-10}\begin{corrige}{10}{Reconnaître une réduction}\demarche{Élever le rapport au cube pour le volume.}
Le facteur des volumes est $0{,}4^3=0{,}064$. Le volume vaut $750\times0{,}064=\res{48\text{ cm}^3}$. Employer 0,4 seul conviendrait aux longueurs.
\label{fig-cor-15-10}\figureGeo[corrige]{transfo-10}
\end{corrige}\end{blocLivre}
\begin{blocLivre}{cor-15-11}\begin{corrige}{11}{Une frise}\demarche{Compter les intervalles entre motifs.}
Il s’agit d’une translation horizontale de 4 cm vers la droite. Entre les motifs 1 et 6, il y a cinq déplacements : $\res{20\text{ cm}}$.
\label{fig-cor-15-11}\figureGeo[corrige]{transfo-11}
\end{corrige}\end{blocLivre}
\begin{blocLivre}{cor-15-12}\begin{corrige}{12}{Une rosace}\demarche{Partager le tour complet en parts égales.}
Un pas vaut $360/8=\res{45^\circ}$. Trois pas valent $\res{135^\circ}$ dans le sens choisi. Huit pas ramènent au départ.
\label{fig-cor-15-12}\figureGeo[corrige]{transfo-12}
\end{corrige}\end{blocLivre}
\begin{blocLivre}{cor-15-13}\begin{corrige}{13}{Retrouver un centre}\demarche{Croiser les droites point-image puis contrôler le rapport.}
Les droites AA’ et BB’ se coupent en $\res{O(-1;-1)}$. On vérifie que A est le milieu de OA’ et B le milieu de OB’ : $(O+A')/2=A$ et $(O+B')/2=B$, coordonnée par coordonnée.
\label{fig-cor-15-13}\figureGeo[corrige]{transfo-13}
\end{corrige}\end{blocLivre}
\begin{blocLivre}{cor-15-14}\begin{corrige}{14}{Une aire donne-t-elle le signe ?}\demarche{Une information d’aire fixe seulement la valeur absolue.}
$k^2=75/12=25/4$. Donc $\res{k=2{,}5\text{ ou }k=-2{,}5}$. L’aire ne permet pas de choisir le signe ; il faut savoir de quel côté du centre sont les sommets homologues.
\label{fig-cor-15-14}\figureGeo[corrige]{transfo-14}
\end{corrige}\end{blocLivre}
\begin{blocLivre}{cor-15-15}\begin{corrige}{15}{Deux transformations successives}\demarche{Suivre séparément côté du centre et distance.}
Le premier changement place le point de l’autre côté à distance $2OA$. Le second change encore de côté et multiplie cette distance par 0,5. Le point revient exactement en A. Le rapport global vaut $\res{(-2)(-0{,}5)=1}$.
\label{fig-cor-15-15}\figureGeo[corrige]{transfo-15}
\end{corrige}\end{blocLivre}
